\keys{Software Engineering; Software Factory; Authentication and Security; Software Quality; Full-Stack Development.}

\begin{abstract}{Full-Stack Evolution and Refactoring of the SIPROS System: Security, Quality, and Interface Improvements in a Software Factory}

\textbf{Context:}
Academic training in Software Engineering requires practical experience to consolidate theoretical learning. Working in environments that simulate the industry, such as an academic Software Factory, provides a realistic transition, demanding adaptation to existing architectures, continuous quality assurance, and the adoption of modern security protocols.

\textbf{Objective:}
This work aims to describe and analyze the practical experience of the full-stack evolution and refactoring of the Integrated Selection Processes System (SIPROS) at UFG, highlighting the implementation of robust security mechanisms, user interface improvements, and the adoption of quality assurance practices.

\textbf{Method:}
The work was developed using action research methodology within the context of the Software Factory course during the 2026/1 semester. Agile methods were utilized for planning and iterative development. The execution involved rigorous Verification and Validation (V\&V) processes between Use Cases and prototypes, as well as the systematic application of peer reviews (Code Review) and automated testing.

\textbf{Results:}
The intervention resulted in the successful integration of Google authentication via Keycloak, adopting the PKCE (Proof Key for Code Exchange) extension to mitigate structural vulnerabilities in the Angular front-end. Simultaneously, the V\&V practices allowed the correction of dozens of architectural and usability inconsistencies, while the formal structuring of the unit test suite ensured the back-end's integrity against regressions.

\textbf{Conclusions:}
The results show that implementing security protocols in Single Page Applications (SPAs) requires architectural rigor that affects the entire system flow. The experience confirmed that the constant employment of quality practices (testing and V\&V) is indispensable in collaborative projects. Finally, working in the Software Factory provided invaluable technical and behavioral maturation, genuinely bringing the academic experience closer to the reality of the technology market.

\end{abstract}
